\section{A question with several parts (12 points)}

The question goes here. Every \texttt{section} is one question: the number
sits in a filled box, the title next to it, and a shadowed rule underneath.

\begin{enumerate}
    \item First part.
    \item Second part, with a footnote.\footnote{Footnote}
    \item Third part.
\end{enumerate}

\section{A question with subsections (10 points)}

When a question splits into parts that each want their own heading, use
subsections.

\subsection{First part}

Subsection numbers get an outlined box rather than a filled one, so the two
levels stay apart.

\subsection{Second part}

Body text is set in Frutiger; headings and numbers in PixeledSans.

\section{A question with a hint box (20 points)}

\begin{neobox}[neotitlefont=\pixelfont]{Hint}
    \texttt{neobox} is the same box that holds the rules at the front. The
    \texttt{neotitlefont} key sets the title font, and the optional argument
    takes any \texttt{tcolorbox} key.
\end{neobox}
